% CMPE 249 HW (Track A) report. IEEE conference format, 6-10 pages including references.
% Build: latexmk -pdf -outdir=build main.tex   (Overleaf and MacTeX ship IEEEtran)
% Tables in tables/ are generated by `python -m det_eval.paper_tables`; figures by `python -m det_eval.run figures`
% and `python -m det_eval.gallery`. Rule for every figure caption: state the takeaway, not the axes.
\documentclass[conference]{IEEEtran}
\usepackage[T1]{fontenc}
\usepackage{graphicx}
\usepackage{booktabs}
\usepackage{amsmath}
\usepackage{xcolor}
\usepackage{url}
\usepackage[hidelinks]{hyperref}

\graphicspath{{figures/}}
\newcommand{\todo}[1]{\textcolor{red}{[TODO: #1]}}
\newcommand{\repourl}{\url{https://github.com/Adithya-Rajendran/HW1-2D-Detection}}
% \figslot{file}{width}{placeholder text}: shows the figure once it has been generated
\newcommand{\figslot}[3]{\IfFileExists{figures/#1}{\includegraphics[width=#2]{#1}}%
  {\fbox{\parbox[c][1.1in][c]{0.95#2}{\centering\footnotesize\textit{#3}}}}}
% \tabin{name}: expandable input of a generated table body (LaTeX's \input is not expandable inside tabular)
\makeatletter\newcommand{\tabin}[1]{\@@input tables/#1.tex }\makeatother

\begin{document}

\title{One Flag and Two Gaps: What Zero-Shot\\COCO-to-Driving Detection Results Hide}
\author{\IEEEauthorblockN{Adithya Rajendran}
\IEEEauthorblockA{San Jos\'e State University\\CMPE 249 Intelligent Autonomous Systems, Fall 2026 (Track A)}}
\maketitle

\begin{abstract}
COCO-pretrained detectors lose half their AP on KITTI. We evaluate seven configurations from four families (YOLO26
with three post-processing variants, RT-DETR, open-vocabulary YOLOE, the Cosmos~3 Edge VLM, and YOLO26 fine-tuned) zero-shot and fine-tuned, with a devkit-faithful
evaluator and paired bootstrap tests, on a Jetson Orin Nano. Three findings change how such comparisons should be
read. First, one post-processing flag decides an architecture comparison: Ultralytics' predictor keeps one class per
box, and enabling multi-label NMS adds 1.6 AP on KITTI and beats the NMS-free head it otherwise trails. Second, the
domain gap is two gaps: doubling resolution helps only cars under 40~px, fine-tuning helps only cars above 25~px
(car AP 50.7 to 68.3). Third, 60 to 70\% of audited ``false positives'' on tiny boxes are unlabeled real objects.
On the robot, a TensorRT detector runs in 16~ms; the VLM takes 17~s and finds 11\% of cyclists.
\end{abstract}

\section{Introduction}
Robots inherit their eyes from datasets they will never see again. A mobile base or a vehicle typically runs a
detector pretrained on COCO \cite{lin2014coco} and deployed, unchanged, on its own camera. KITTI \cite{geiger2012kitti}
is the classic test of that transfer: the objects are familiar, but the viewpoint, the scale and the labeling
conventions are not. The course tutorial \cite{liu2026ngdet} reports zero-shot KITTI mAP of 0.23 to 0.28 for COCO
models; we reproduce that range on the full validation split and ask what it is made of, and what a robot can do
about it on an 8~GB Jetson Orin Nano.

The answer turns out to depend as much on evaluation choices as on models, so we organize the paper around three
findings rather than a leaderboard:
\begin{itemize}
  \item \textbf{F1, one flag.} The NMS-free head's KITTI advantage over the NMS head (+0.7 AP, +11 AR) is multi-label
  output, not the absence of NMS. One predictor flag reverses the comparison and explains the gap to published COCO
  numbers (Section~\ref{sec:flag}).
  \item \textbf{F2, two gaps.} Resolution recovers far-field objects and nothing else; in-domain fine-tuning recovers
  near-field objects and nothing else (Section~\ref{sec:gaps}).
  \item \textbf{F3, missing labels.} Most of the small ``false positives'' that separate models are real objects
  KITTI never labeled (Section~\ref{sec:failures}).
\end{itemize}
Around them we report the full Track~A comparison: closed-set vs.\ open-vocabulary vs.\ VLM, speed vs.\ accuracy on
the Jetson, zero-shot vs.\ fine-tuned, and a typed failure analysis. Section~\ref{sec:practice} turns the findings
into a checklist for anyone comparing detectors across datasets, and the harness, configs, manifests and logs are
public (\repourl).

\section{Background and Related Work}
\textbf{Detector families.} Two-stage detectors \cite{ren2015faster} propose regions and then classify them; dense
one-stage detectors \cite{redmon2016yolo} predict a box at every feature-map location; query-based detectors
\cite{carion2020detr,zhu2021deformable,zhao2024rtdetr} predict a fixed set. The axis that separates them is target
assignment: IoU and task-aligned matching give an object many positive anchors (one-to-many), Hungarian matching
gives it one (one-to-one).

\textbf{From NMS to end-to-end detection.} One-to-many training makes duplicates at inference, so dense detectors
need non-maximum suppression (NMS). DETR's bipartite matching removes it \cite{carion2020detr}; YOLOv10 trains a
one-to-one head alongside a one-to-many head \cite{wang2024yolov10}, and YOLO26 ships both heads on the same weights
\cite{yolo26docs}, which lets us compare them with everything else fixed. RT-DETR argues that NMS makes latency
scene-dependent \cite{zhao2024rtdetr}. Low-score boxes also matter downstream: ByteTrack recovers occluded objects
by associating detections far below the usual threshold \cite{zhang2022bytetrack}.

\textbf{Open vocabulary and VLMs.} Grounding DINO \cite{liu2024groundingdino} and YOLOE \cite{wang2025yoloe} take
class names as text prompts; general VLMs now emit boxes as text \cite{qwen3vl}. Cosmos~3 \cite{nvidia2026cosmos3}
is NVIDIA's physical-AI family, and its 2.4B Edge variant runs on the Orin Nano \cite{cosmos3edge_jetson}.

\textbf{Diagnosing errors.} Hoiem et al.\ \cite{hoiem2012diagnosing} and TIDE \cite{bolya2020tide} decompose
detector errors into localization, classification, duplicate, background and missed types; our error census adapts
their types to KITTI's ignore regions and adds a manual audit of the ``background'' type.

\section{Experimental Setup}
\subsection{Models}
Table~\ref{tab:models} lists the configurations: three closed-set COCO detectors from two families, one
open-vocabulary detector, one VLM, and YOLO26m fine-tuned on KITTI. YOLO26m and RT-DETR-l have nearly identical
published COCO AP, so a difference on KITTI cannot be explained by in-domain accuracy.

\begin{table}[t]
\centering
\caption{Model configurations. Published accuracy is COCO val AP, or LVIS minival AP with text prompts for YOLOE.
``Multi-label'' runs the same NMS head with \texttt{multi\_label=True}, as Ultralytics' validator does.}
\label{tab:models}
\footnotesize
\setlength{\tabcolsep}{4pt}
\begin{tabular}{@{}lll@{}}
\toprule
Config & Family & Published \\
\midrule
YOLO26m, NMS & dense CNN, one-to-many & 53.1 \\
YOLO26m, NMS multi-label & same head, multi-label NMS & -- \\
YOLO26m, NMS-free & same weights, one-to-one & 52.5 \\
RT-DETR-l & queries, Hungarian matching & 53.0 \\
YOLOE-26m & open-vocab, text prompts & 35.4 (LVIS) \\
Cosmos~3 Edge & VLM (2.4B), 2D grounding & n/a \\
YOLO26m-FT & fine-tuned on KITTI train & n/a \\
\bottomrule
\end{tabular}
\end{table}

\subsection{Datasets and taxonomy mapping}
COCO val2017 (5{,}000 images) is the in-domain reference. KITTI 2D object detection \cite{geiger2012kitti} is
evaluated on the standard 3{,}769-frame validation half of the labeled set \cite{chen2015_3dop}: 14{,}385 cars,
2{,}280 pedestrians, 893 cyclists, 606 trucks and 287 trams. Fine-tuning uses the other 3{,}712 frames only. The
VLM comparison uses a fixed random 300-frame subset (seed 0). As a second driving domain we use 3{,}000 random
key frames (seed 0) of nuImages val \cite{caesar2020nuscenes}, which the tutorial recommends as the reliable 2D
driving benchmark. Table~\ref{tab:mapping} gives the KITTI mapping; the nuImages mapping follows the same rules and
is in the repository.

\begin{table}[t]
\centering
\caption{KITTI-to-COCO mapping. ``Ignore'' boxes are neither true positives nor false positives.}
\label{tab:mapping}
\footnotesize
\begin{tabular}{@{}p{2.3cm}p{5.5cm}@{}}
\toprule
KITTI & Treatment \\
\midrule
Car, Pedestrian, Truck, Tram & COCO car, person, truck, train \\
Van, Person\_sitting & ignored for car / person, IoU-matched (devkit neighboring classes) \\
Cyclist & ignored for person, matched by containment (a COCO person box covers only the rider); a real class for open-vocabulary models and the fine-tuned model \\
DontCare & ignore region for all classes (intersection over detection area) \\
Misc & dropped \\
\bottomrule
\end{tabular}
\end{table}

\subsection{Evaluation protocol}
Every model is scored by one pycocotools-based evaluator \cite{cocoapi} that applies the ignore rules. AP runs keep
boxes down to confidence 0.001 and 100 per image; the deployment operating point is confidence 0.25. We report
COCO metrics (AP, AP$_{50}$, AP$_{75}$, AP$_{S/M/L}$, AR$_{100}$, per-class AP), AP by box height in pixels, and
KITTI-style AP$|_{R40}$ at the devkit IoU (0.7 for vehicles, 0.5 for pedestrians and cyclists). Intervals are 95\%
image-level bootstrap intervals \cite{efron1993bootstrap}; differences between two models use a paired bootstrap
(300 to 500 resamples, the same image weights applied to both), which our tests show reproduces pycocotools exactly
at unit weights. The VLM emits no scores, so it is compared by precision, recall and F1 at IoU 0.5 and 0.7.
Closed-set models get a rule-based cyclist (a person box over a bicycle box). Latency is measured at batch 1 on 500
KITTI frames after 50 warm-up frames, at confidence 0.25, synchronized, from decoded image to boxes; peak memory is
the PyTorch allocator peak (unified memory has no per-process counter).

\subsection{Hardware, software and reproducibility}
All runs use a Jetson Orin Nano Super 8~GB (JetPack 6.2.3, L4T R36.5.2) in MAXN\_SUPER mode with clocks locked by
\texttt{jetson\_clocks} (GPU 1.02~GHz, EMC 3.2~GHz), torch 2.10, Ultralytics 8.4.183 \cite{ultralytics} and
TensorRT 10.3 FP16 engines; Cosmos~3 Edge runs in BF16 through Hugging Face Transformers. Inference is FP16; on
1{,}000 COCO images FP32 gives 54.28 AP and FP16 54.16, so precision is not a factor. Fine-tuning ran in the
Ultralytics JetPack~6 container (same Ultralytics and torch versions). Every run writes a manifest with its command,
settings and versions; Appendix~\ref{app:repro} lists them.

\section{Results}
\subsection{Baseline and reproduction}
Table~\ref{tab:main} gives the main results. COCO AP of 52 to 53 falls to 26 to 28 on KITTI for every closed-set
model; YOLOE falls to 18. RT-DETR-l reproduces its published COCO number to within 0.3 AP and the NMS-free YOLO26m
head to within 0.3, but the NMS head scores 52.2 against a published 53.1. That gap is the first clue to F1.

\textit{Course harness.} We also ran the tutorial's own \texttt{run\_eval} \cite{liu2026ngdet} unmodified on the
Orin Nano (Table~\ref{tab:ngdet}). Its exact command gives 23.6 to 24.0 mAP against the published 26.2 to 27.9;
lowering the score threshold to 0.001 and using 1,000 frames instead of 120 each move the numbers by under one point,
so neither explains the remaining gap. Three protocol switches, applied to a separate patched copy that matches the
original exactly when they are off, each change the result by about as much as the gaps between models. Treating
KITTI DontCare regions as ignore regions, as the devkit does, adds 0.5 to 1.1 mAP to every model. Splitting the
merged \emph{vehicle} class into car and truck (\texttt{driving5}) is what makes multi-label NMS matter: it adds 0.6
mAP and 7.7 AR there and nothing under \texttt{driving3}, where a car/truck duplicate is the same class. The switches
reorder the models: RT-DETR-R50 narrowly leads under the tutorial's protocol, YOLO26x with multi-label NMS under the
five-class one. Under both taxonomies cyclist AP is 0.0 for every model, because a COCO \emph{bicycle} box never
matches a KITTI \emph{Cyclist} box drawn around rider and bicycle; a third of the \texttt{driving3} mean is fixed at
zero, which our rider-plus-bicycle composition (Table~\ref{tab:mapping}) avoids.

\begin{table}[t]
\centering
\caption{The course harness on KITTI, Orin Nano (mAP in \%, 1,000 frames unless noted). ML: multi-label NMS;
NF: NMS-free head. Each protocol switch moves every model by about as much as the gaps between models.}
\label{tab:ngdet}
\scriptsize
\setlength{\tabcolsep}{3pt}
\begin{tabular}{@{}lccccc@{}}
\toprule
Protocol & \multicolumn{3}{c}{YOLO26x} & RT-DETR & F.\ R-CNN \\
 & NMS & ML & NF & R50 & v2 \\
\midrule
\tabin{ngdet}
\bottomrule
\end{tabular}
\end{table}

\begin{table*}[t]
\centering
\caption{Main results (AP in \%, 95\% bootstrap interval in brackets; S/M/L are COCO box-area bins). COCO AP halves
on KITTI for every closed-set model; 1280~px input adds 1 to 3 AP. Fine-tuned rows use the same 4-class KITTI ground
truth as the zero-shot rows.}
\label{tab:main}
\scriptsize
\begin{tabular}{@{}lllcccccccc@{}}
\toprule
Model & Data & Input & AP & AP$_{50}$ & AP$_{75}$ & AP$_S$ & AP$_M$ & AP$_L$ & AR$_{100}$ & KITTI mod.\ car / ped \\
\midrule
\tabin{main}
\bottomrule
\end{tabular}
\end{table*}

\subsection{F1: one flag decides the head comparison}\label{sec:flag}
On COCO the NMS and NMS-free heads tie (52.2 vs.\ 52.2). On KITTI the NMS-free head is ahead by a small but
significant 0.7 AP and 11.0 points of class-averaged AR$_{100}$ (Table~\ref{tab:paired}). We expected the reason to
be NMS deleting overlapping parked cars. It is not: per class, recall with every box kept rises by 30 points for
trucks (0.61 to 0.90) and 17 for trams, but only 1.8 for cars, and the error census shows no fewer ``crowded''
misses. At the 0.25 deployment threshold the heads are indistinguishable (recall 0.812 vs.\ 0.806).

The mechanism is in post-processing. The NMS-free head takes its top-$k$ over (anchor, class) pairs, so one anchor
can be reported as both ``car'' and ``truck''. Ultralytics' \texttt{predict()} runs NMS on the arg-max class only,
while its validator, which produced the published numbers, runs NMS with \texttt{multi\_label=True}. Enabling that
one flag in the predictor:
\begin{itemize}
  \item lifts truck recall from 0.61 to 0.90 and tram recall from 0.46 to 0.66, matching the NMS-free head;
  \item adds 1.6 AP and 9.5 AR on KITTI, and beats the NMS-free head by 0.9 AP (paired interval $-$1.15 to $-$0.57
  for NMS-free minus multi-label) while trailing it by 1.6 AR;
  \item raises our COCO number from 52.2 to 52.8. Ultralytics' own validator gives 53.0 (NMS) and 52.5 (NMS-free) on
  the same weights, matching the published 53.1 and 52.5; the last 0.2 is batching and letterboxing.
\end{itemize}
The flag matters more on KITTI than on COCO because KITTI's vans and box trucks sit on COCO's car/truck boundary;
the NMS-free head already leads by 3.3 truck AP on COCO itself (Fig.~\ref{fig:nmsgap}). The NMS-free head keeps a
real advantage, latency: under TensorRT it runs at 15.6~ms p50 against 18.0~ms (Table~\ref{tab:latency}), almost all
of it the NMS step. Neither head's latency depends on scene clutter with clocks locked (Fig.~\ref{fig:latency}).
\todo{nuImages: does the multi-label effect replicate on car/truck/bus?}

\begin{figure}[t]
\centering
\figslot{f3_nms_gap.pdf}{\columnwidth}{F3: AP(NMS head) minus AP(NMS-free head), COCO vs.\ KITTI}
\caption{The head gap is a class-boundary effect present in domain. The heads differ by under 0.6 AP on cars and
people in both domains; the NMS-free head's truck advantage already exists on COCO (3.3 AP) and shrinks on KITTI
(1.2); only trams grow (0.3 to 2.2).}
\label{fig:nmsgap}
\end{figure}

\begin{table*}[t]
\centering
\caption{Paired differences (points, 95\% paired bootstrap interval; $^{*}$ excludes zero). AR is class-averaged
AR$_{100}$; ``AP $<$25'' is the class-averaged AP on boxes under 25~px tall.}
\label{tab:paired}
\footnotesize
\begin{tabular}{@{}llccc@{}}
\toprule
Comparison & Data & $\Delta$AP & $\Delta$AR & $\Delta$AP $<$25 \\
\midrule
\tabin{paired}
\bottomrule
\end{tabular}
\end{table*}

\subsection{F2: the domain gap is two gaps}\label{sec:gaps}
Doubling the input from 640 to 1280~px raises KITTI AP by 2.4 for the NMS head (paired interval +1.8 to +3.1).
Fig.~\ref{fig:height} shows where: on cars under 25~px tall, YOLO26m gains 7.5 AP (44.8 to 52.3) and YOLOE 9.6; on
25 to 40~px cars about 3; above 40~px nothing (59.6 vs.\ 58.9 above 80~px). With KITTI's focal length of about
720~px, a car 5~ft tall is 25~px tall at roughly 140~ft, so resolution is a far-field fix. We plot cars because
they are the only class with enough small objects: the split has 32 pedestrians and 10 trams under 25~px, and
class-averaged height-bin AP is dominated by that noise.

The near field is the larger loss, and it is localization. On COCO YOLO26m keeps 75\% of its AP$_{50}$ at
AP$_{50:95}$ (52.2 / 69.5); on KITTI it keeps 52\% (26.2 / 50.1). At the deployment threshold, 48\% of its misses
are ``loose'' (a box on the object at IoU below 0.5): large near cars, often truncated at the border, get covered
by partial boxes (Fig.~\ref{fig:gallery}a). Fine-tuning removes exactly this part. At 640~px, car AP rises from 59.6
to 81.1 above 80~px and from 48.9 to 68.9 at 40 to 80~px, loose misses on the subset fall from 149 to 13, but cars
under 25~px do not move (44.8 vs.\ 45.6). In-domain data fixes the near field; only more pixels fix the far field.

\begin{figure}[t]
\centering
\figslot{f2_height_bins.pdf}{\columnwidth}{F2: car AP by box height}
\caption{Resolution and in-domain data fix different halves of the drop. Doubling the input adds about 7 AP on
cars under 25~px and nothing above 40~px; fine-tuning at 640~px (green) adds 20 AP on near cars and nothing under
25~px. Car AP, KITTI val.}
\label{fig:height}
\end{figure}

\subsection{Fine-tuning: what a night on the robot buys}
YOLO26m was fine-tuned from the COCO checkpoint on 3{,}412 KITTI training frames (300 more held out for checkpoint
selection) with five classes named as in COCO (car, person, cyclist, truck, train), which lets Ultralytics copy the
pretrained classifier rows for four of them. Training ran for a fixed 7-hour budget (84 epochs) at 640~px, batch 8,
on the Orin Nano itself. YOLO training has no ignore regions, so Van, Person\_sitting, Misc and DontCare are
background; evaluation still uses the devkit rules. Ultralytics' score on the held-out training frames reached
0.79 AP, but those frames share drives with the training frames; Table~\ref{tab:finetune} uses the validation split.

The result is mixed in an informative way (Table~\ref{tab:finetune}). Cars, with 13{,}195 training boxes, improve
from 50.7 to 68.3 AP and from 74.7 to 90.6 KITTI moderate AP$|_{R40}$, and subset F1 rises from 0.73 to 0.84 as
precision jumps from 0.68 to 0.92. Trucks and trams, with 450 and 206 training boxes, regress (11.3 to 8.0 and 18.5
to 10.3 AP), the forgetting the tutorial also reports for small per-domain sets \cite{liu2026ngdet}. Pedestrians gain
2.6 COCO-style AP but lose 7 KITTI moderate AP at IoU 0.5, and AR drops because the fine-tuned model emits far fewer
low-score boxes. The class-averaged gain is only +2.2 AP (paired interval +0.9 to +3.3), which hides a large car gain
behind two rare-class regressions. We hypothesized that training on Vans as background teaches the model not to fire
on van-shaped trucks; an ablation that grays out Van and Person\_sitting regions during training tests this.
\todo{maskvan result}

\begin{table}[t]
\centering
\caption{Per-class AP on KITTI val at 640~px. Zero-shot cyclist AP marked $^{r}$ uses the person-over-bicycle rule.
Fine-tuning gains 18 AP on cars, while trucks and trams, with few training boxes, regress.}
\label{tab:finetune}
\scriptsize
\setlength{\tabcolsep}{1.8pt}
\begin{tabular}{@{}lccccccc@{}}
\toprule
Model & AP & car & person & truck & train & cyclist & AR$_{100}$ \\
\midrule
\tabin{finetune}
\bottomrule
\end{tabular}
\end{table}

\subsection{Closed-set, open-vocabulary and VLM}
Table~\ref{tab:vlm} compares all families on the 300-frame subset. Cosmos~3 Edge parses cleanly (299 of 300 replies
valid JSON, one salvaged) and finds 71\% of objects, but at 0.61 precision its F1 is 0.66, against 0.73 for YOLO26m
and 0.84 fine-tuned, and its boxes are looser (F1 at IoU 0.7: 0.50 vs.\ 0.57). It takes 17.3~s per frame (p90 36~s),
about 1{,}100 times the TensorRT detector (Fig.~\ref{fig:speedacc}). YOLOE-26m is the most precise zero-shot model
(0.78) but the least complete (recall 0.61), and trails the closed-set YOLO26m by 7.8 AP on KITTI.

No method covers cyclists, the class COCO lacks. Cosmos~3 Edge finds 11\%, mostly because it boxes the rider as a
pedestrian (Fig.~\ref{fig:gallery}e); YOLOE, prompted with ``cyclist'', finds none above 0.25. The zero-cost rule on a
closed-set detector reaches 24 to 33\% recall, and 681 labeled training cyclists move the fine-tuned model to only
26\% at the deployment threshold, though its cyclist AP (18.9) is the best of any model.

\begin{figure}[t]
\centering
\figslot{f5_speed_accuracy.pdf}{\columnwidth}{F5: F1 at IoU 0.5 vs.\ latency per frame (log scale)}
\caption{The VLM sits three orders of magnitude to the right of the detectors and below them in F1; open
vocabulary does not buy coverage of the missing class. Detectors at conf 0.25, TensorRT FP16.}
\label{fig:speedacc}
\end{figure}

\begin{table}[t]
\centering
\caption{300-frame subset. Detectors at conf 0.25, the VLM with every box. Cyclist recall at the same operating
point; closed-set detectors use the person-over-bicycle rule. Latency is TensorRT p50 (ms) for detectors and
end-to-end generation for the VLM.}
\label{tab:vlm}
\scriptsize
\setlength{\tabcolsep}{2.5pt}
\begin{tabular}{@{}lccccccc@{}}
\toprule
& \multicolumn{3}{c}{IoU 0.5} & IoU 0.7 & \multicolumn{2}{c}{Cyclist} & \\
\cmidrule(lr){2-4}\cmidrule(lr){5-5}\cmidrule(lr){6-7}
Model & P & R & F1 & F1 & AP$_{50}$ & R & ms \\
\midrule
\tabin{vlm300}
\bottomrule
\end{tabular}
\end{table}

\subsection{Speed and memory on the robot}
TensorRT FP16 halves YOLO latency and cuts RT-DETR's by almost 3$\times$ (Table~\ref{tab:latency}). Every detector
fits in under 150~MB of allocator memory, so the Orin Nano's 8~GB is not the constraint for detection; it is for the
VLM, which needs the whole board. RT-DETR-l is 0.6 to 1.3 AP ahead of YOLO26m's default heads but 0.3 behind the
multi-label NMS head, at twice the latency.

\begin{figure}[t]
\centering
\figslot{f4_latency_vs_clutter.pdf}{\columnwidth}{F4: p50 and p99 latency vs.\ objects per frame}
\caption{Under TensorRT with locked clocks, latency does not move with scene clutter for any model, at p50 or p99;
the NMS-free head is the fastest at every clutter level.}
\label{fig:latency}
\end{figure}

\begin{table}[t]
\centering
\caption{Latency on the Jetson Orin Nano (ms, batch 1, 640~px, conf 0.25, 500 KITTI frames) and peak allocator
memory. The NMS-free head saves 1.9~ms of post-processing.}
\label{tab:latency}
\footnotesize
\setlength{\tabcolsep}{3.5pt}
\begin{tabular}{@{}lcccccc@{}}
\toprule
& \multicolumn{2}{c}{PyTorch FP16} & \multicolumn{3}{c}{TensorRT FP16} & Peak \\
\cmidrule(lr){2-3}\cmidrule(lr){4-6}
Model & p50 & p99 & p50 & p99 & post & MB \\
\midrule
\tabin{latency}
\bottomrule
\end{tabular}
\end{table}

\subsection{Second domain: nuImages}
\todo{nuImages results: zero-shot AP for all detectors, multi-label effect on car/truck/bus, 640 vs.\ 1280 by height,
and whether KITTI fine-tuning transfers or forgets.}

\section{Failure Analysis and F3: what the labels miss}\label{sec:failures}
Every error at the deployment threshold gets one type (Table~\ref{tab:failures}; definitions follow
\cite{hoiem2012diagnosing,bolya2020tide} and are in \texttt{det\_eval/failures.py}). Loose localization is the largest
category for every zero-shot model, in both misses and false positives, and the one fine-tuning removes. RT-DETR-l's
higher recall comes with nearly three times YOLO26m's false positives at the same threshold because its scores are
calibrated lower: a shared threshold means something different for each model.

The ``tiny, unlabeled'' type deserves its own finding. KITTI marks some distant regions DontCare but leaves many
distant objects unlabeled. In random samples of 30 such false positives, 21 of YOLO26m's (70\%, Wilson interval 52
to 83\%) and 18 of Cosmos's (60\%, 42 to 75\%) were real unlabeled vehicles or people (Fig.~\ref{fig:gallery}d).
Cosmos has 246 of these on the subset against YOLO26m's 79, so the metric penalizes it most: its true precision is
closer to 0.70 than 0.61. The genuine hallucinations differ by model: YOLO26m calls
roadside delineator posts and barriers ``person'' (4 of its 8; Fig.~\ref{fig:gallery}c), a phantom-braking risk;
Cosmos boxes signs, bollards and parts of labeled cars. Taxonomy drives the ``confused'' type: COCO calls a delivery
van a truck, KITTI calls it a Van, and the devkit ignores Van only for car (Fig.~\ref{fig:gallery}b).

\begin{table*}[t]
\centering
\caption{Error census on the 300-frame subset at the deployment operating point (IoU 0.5). Loose localization
dominates every zero-shot model and nearly vanishes after fine-tuning; RT-DETR-l trades nearly three times the false
positives for its recall; most of Cosmos~3 Edge's false positives are tiny boxes, of which an audited 60\% are real
unlabeled objects.}
\label{tab:failures}
\scriptsize
\setlength{\tabcolsep}{3.2pt}
\begin{tabular}{@{}lc|cccccc|ccccc@{}}
\toprule
& & \multicolumn{6}{c|}{Missed objects} & \multicolumn{5}{c}{False positives} \\
Model & TP & loose & crowded & confused & small & occluded & missed & loose & duplicate & confused & tiny, unlabeled & other \\
\midrule
\tabin{failures}
\bottomrule
\end{tabular}
\end{table*}

\begin{figure*}[t]
\centering
\figslot{f6_gallery.pdf}{\textwidth}{F6: failure gallery}
\caption{One example per failure type. White: KITTI ground truth; dashed blue: DontCare; blue: YOLO26m (NMS) boxes
at conf 0.25; violet: Cosmos~3 Edge; red: the object or box the panel is about. (a)~Loose: 48\% of YOLO26m's misses
on KITTI val. (b)~Taxonomy: COCO ``truck'' on a KITTI Van. (c)~A delineator post called a person. (d)~A real
distant car with no label, scored as a false positive. (e,f)~Cosmos~3 Edge boxes riders, not cyclists, and skips a
large near car.}
\label{fig:gallery}
\end{figure*}

\section{Recommendations for Cross-Dataset Detection Comparisons}\label{sec:practice}
\begin{enumerate}
  \item \textbf{Fix and report the post-processing.} State whether NMS is multi-label; with Ultralytics, compare
  heads under the same setting or the comparison measures the flag, not the head.
  \item \textbf{Reproduce in-domain first.} A 0.9 AP COCO gap was the clue to item 1.
  \item \textbf{Score with the target dataset's ignore rules.} Without DontCare and neighbor classes, KITTI
  charges every model for objects it was never asked to find.
  \item \textbf{Report per-class and per-height results, and test differences with a paired bootstrap.} Class means
  hide an 18-point car gain behind two rare-class regressions, and class-averaged small-object AP can rest on ten
  objects.
  \item \textbf{Audit the background errors before ranking.} Sixty to seventy percent of tiny false positives here
  were real objects.
  \item \textbf{Compare at matched operating points.} A shared confidence threshold is not a shared precision.
\end{enumerate}

\section{Limitations}
One board, and the VLM is scored without confidence on 300 frames with one prompt. Mapping choices matter: Van is
ignored only for car. The fine-tuned model is one 7-hour run with default hyperparameters and no ignore regions,
a baseline rather than a tuned result. The label audit covers 60 boxes and quantifies, but does not correct, KITTI's
missing labels. Latency is measured with locked clocks and may differ under thermal throttling.

\section{Conclusion}
Zero-shot COCO-to-driving numbers hide one flag and two gaps. The NMS-free head's KITTI advantage is multi-label
output, which an NMS head gets from one setting, and its lasting benefit is 2.4~ms of latency. Resolution recovers
far-field cars and nothing else; a night of fine-tuning on the robot recovers near-field cars (50.7 to 68.3 AP) and
nothing else. A 2.4B VLM runs on the robot but trails a 20M-parameter detector by 0.07 F1 and three orders of
magnitude in speed, and no method we tried made cyclists reliable.

\section*{AI Tool Use}
As the assignment allows, Claude (Anthropic) assisted with writing the evaluation harness, debugging the Jetson
setup and drafting this report. Every number in the paper is generated from the run outputs by
\texttt{det\_eval.paper\_tables}, the evaluator is covered by unit tests against pycocotools, and the error audit
was judged by inspecting each crop.

\appendix
\section*{Reproducibility}\label{app:repro}
Code, configs (\texttt{configs/experiments.yaml} holds every setting), run manifests and logs: \repourl. The
pipeline is: \texttt{python -m det\_eval.run gt}; \texttt{predict}; \texttt{evaluate}; \texttt{latency [--engine]};
\texttt{vlm --model cosmos3\_edge}; \texttt{report}; \texttt{figures}. Fine-tuning:
\texttt{python -m det\_eval.finetune split}, \texttt{export}, \texttt{train --hours 7}. Paired tests:
\texttt{python -m det\_eval.compare --gt \dots\ --a \dots\ --b \dots}. Table~\ref{tab:repro} lists every prediction run
and Table~\ref{tab:commands} every command recorded in the manifests.

\begin{table*}[h]
\centering
\caption{Every prediction run, from its manifest.}
\label{tab:repro}
\scriptsize
\setlength{\tabcolsep}{3pt}
\begin{tabular}{@{}llllllll@{}}
\toprule
Run & Checkpoint & Data (frames) & Input & conf / head & Precision & Hardware & Framework \\
\midrule
\tabin{repro}
\bottomrule
\end{tabular}
\end{table*}

\begin{table*}[h]
\centering
\caption{Commands recorded in the run manifests.}
\label{tab:commands}
\scriptsize
\begin{tabular}{@{}l@{}}
\toprule
\tabin{commands}
\bottomrule
\end{tabular}
\end{table*}

\bibliographystyle{IEEEtran}
\bibliography{refs}

\end{document}
